% 01-problem.tex: the problem Polaris is trying to solve.
\section{The problem}
\label{sec:problem}

A person in the United States carries six to eight credentials that do not talk to each other:
a driver's licence, a passport, a Social Security card, a Real ID, a voter registration, an
insurance card. Each is a separate artifact, signed by a separate authority to a separate standard,
with no shared revocation path and no shared audit trail. Consolidating them into one active
credential record per person, verified through context-scoped events, is the easy half of the
problem, and it is still the problem's shape.

The hard half is what happens when the system works. An identity system that succeeds becomes the
doorway to everything, and a doorway is a place where power concentrates. Three failure modes
follow, and every layer of Polaris exists to answer one of them.

\begin{description}
  \item[The system lies about the past.] An operator, or an insider with database access,
    revokes a token and later erases the record; issues a credential under procedures the public
    was never told of; or shows one history to an auditor and another to a court. A national
    identity system cannot be allowed to retroactively rewrite what it did.
  \item[The system watches.] Every verification is a fact about where a person was and what they
    were doing. A log of those facts is a surveillance backbone whether or not anyone intended one,
    and it stays one after the operator's intentions change.
  \item[The system coerces.] A holder can be compelled: a border officer, an employer, or a
    thief can force a credential to be presented, and a doorway that every service requires can
    force a population to enrol. An identity system is also the most efficient instrument of
    exclusion ever built, if the list of who is not in it can be produced on demand.
\end{description}

The conventional answer to all three is policy: rules about retention, access, audit and consent,
enforced by the people who run the system. Polaris's constitution (\code{MISSION.md}) takes the
position that policy is worthless as the only thing holding a promise up, because the people who
run the system are exactly the ones the promise is meant to bind. The promise has to be enforced
where it cannot be talked around: in the schema, in triggers, in unique indexes, in check
constraints, in signatures that anyone can verify without asking the system, and in logs that other
people watch. The application is not the security boundary. The database is, and the layers around
the database exist to make even the database's claims checkable by strangers.

\subsection{Why post-quantum from the first day, and what the phrase can claim}

Identity data has a secrecy horizon measured in generations, and migrating a national system takes
a decade. By Mosca's inequality, which compares the years data must stay secret plus the years a
migration takes against the years until a quantum computer can break today's signatures, a credential
system deployed on classical cryptography in 2026 is a deferred compromise rather than a working
system (Appendix~\ref{app:from-v1}). At
\code{1.0.0-rc.7} the default signing algorithm is ML-DSA-65 (FIPS 204), ML-DSA-87 is accepted beside
it, the zero-knowledge proof system is hash-based with no trusted setup, and the public TLS edge and
one of the two internal hops negotiate a hybrid post-quantum key exchange.

\emph{Post-quantum} is accurate about the algorithms. It is not a claim the
repository can make about security, because ML-DSA's security rests on lattice assumptions that could
fall to mathematics rather than to hardware, and because the default development path signs with a
labelled placeholder rather than a lattice signature at all. The defensible claim is \emph{algorithm
agility under an audited migration path}: the algorithm is a row in the schema and never a constant,
a population can be re-signed under a new one on a path that runs on every push, and what a break
would cost is written down. Section~\ref{sec:signatures} gives the argument and the cost; an
invariant check refuses any outward surface that says more.

\subsection{The progression this paper follows}

Each of the following sentences names a layer, and each is true of the layer before it.

\begin{itemize}
  \item A credential's authenticity is not its current authorization. A signature stays genuine
    forever; the authority under it can be withdrawn this afternoon.
  \item Current authorization is not privacy. A verifier can learn that a credential is valid and
    still learn far more than it needed to.
  \item Privacy is not freedom from coercion. A system can know almost nothing about a person and
    still be the doorway every service demands.
  \item Federation is not institutional independence. Two instances of the same code run by the
    same author prove that the protocol works, not that independent institutions run it.
  \item A transparency log is not protection from a split view. One observer can be shown a private
    fork of history and never know.
  \item A green test is not proof that the test asked the right question.
  \item A system that has checked itself is not a system that anyone else has checked. Every
    instrument in the repository passed over a defect that the first outside implementation refused
    on sight.
\end{itemize}

The last two sentences are the ones the project has learned most expensively. Section~\ref{sec:senses}
is about the sixth, and Section~\ref{sec:verifiers} carries the seventh.
